
\section{Related Work}
\subsection{LLM-based Autonomous Agents}
LLM-based autonomous agents are an advanced AI system using a core LLM to manage external functional modules and interact with the world~\citep{ding2024understanding}.  
%In the early stage, LLM-based agents possess very simple architectures, which just directly query LLMs to conduct decision-making and take action~\citep{huang2022language,ahn2022can}. 
Recent studies have equipped LLM agents with several LLM-centric functional modules including planning~\citep{hao2023reasoning,zeng2024perceive,shao2025division}, reasoning~\citep{wei2022chain,yao2024tree,shang2024defint,xu2025towards}, using tools~\citep{shen2024hugginggpt,schick2024toolformer}, and monitoring memory~\citep{wang2024voyager,park2023generative}, greatly enhancing the capabilities of LLM agents. 
Along with the improvement of the single agent, there's another line of work trying to build more advanced multi-agent systems by strategically organizing individual agents for both simulation~\citep{li2023metaagents,chen2023agentverse} and targeted task solving~\citep{qian2024chatdev,chen2024large,li2024limp}.
The emergence of more and more sophisticated agent produces remarkable performance improvement, however, their architectures and codebases differ greatly with each other.
The lack of a unified design space and consistent terminologies across individual works makes it hard to compare different agents, understand their evolution routes, and guide new agent design directions.



\subsection{Automatic Design of LLM-based Agents}
LLM-based agent system, as the most advanced AI system, has not yet formed a unified design space and an automatic design approach.
%Existing efforts on automatic building of LLM agents focus on three aspects: engineering-oriented open resources, conceptual framework of LLM agent design, and automatic design of partial components of LLM agents.
Engineering-oriented open resources like LangChain\footnote{https://github.com/langchain-ai/langchain} and BabyAGI\footnote{https://github.com/yoheinakajima/babyagi} have provided convenient ways to build an LLM-centric agentic system, however, they still need human participation to organize different modules and can't support the optimization of the designed agent.
Besides, there have been some conceptual frameworks trying to provide a unified design principle of LLM agents, such as CoALA~\citep{sumers2023cognitive}. 
%It points out four basic modules of language agents from the perspective of cognitive science, including long-term memory, external grounding, internal actions and decision-making.
However, it's still a vision of how LLM agents should be in the future, without providing a practical design framework.
More importantly, there are several recent works that explore the problem of automating (at least part of) the design of LLM agent systems defined on different search spaces. OPRO~\citep{yang2024large} and Promptbreeder~\citep{fernando2024promptbreeder} can be considered as using LLMs to optimize LLM agent defined on prompt space. More relevantly, ADAS~\citep{hu2024automated} proposes to search the entire agentic system defined on code space, enabling the search for LLM agents with more flexible prompts, tool uses, control flows and more.


% Nevertheless, these works over-simplify the structure of agentic systems to pure prompts or codes, unable to provide a clear design space of LLM agents.
% Besides, another common shortcoming of existing works is that all of them are proposed without much consideration of existing designed agent architectures, thus hard to be representative and universal.
% Differently, our proposed modularized LLM agent design space can well cover existing efforts of LLM agent design, and provide a practical approach for automatic agent architecture optimization.